\documentclass[11pt,letterpaper]{article}

\usepackage[spanish,es-noshorthands]{babel}
\usepackage{fontspec}
\setmainfont{Source Serif Pro}
\usepackage{microtype}
\usepackage{geometry}
\usepackage{enumitem}
\usepackage{xcolor}
\usepackage{csquotes}
\usepackage{float}
\usepackage{tikz}
\usetikzlibrary{arrows.meta,positioning,fit,backgrounds}
\usepackage{xurl}
\usepackage[hidelinks]{hyperref}
\usepackage[backend=biber,style=apa,sorting=nyt,maxcitenames=2]{biblatex}
\DeclareLanguageMapping{spanish}{spanish-apa}

\geometry{margin=2.3cm}
\hbadness=10000
\setlength{\emergencystretch}{1em}
\setlength{\parindent}{0pt}
\setlength{\parskip}{0.5em}
\setlist{nosep}
\addbibresource{referencias.bib}
\AtBeginBibliography{\sloppy\small}
\setlength{\bibitemsep}{0.3em}

% Colores: verde para los proyectos que hoy son parte del núcleo del negocio,
% terracota para el que se implementó sin llegar al núcleo, gris para los
% pilotos que no avanzaron y azul oscuro para los rótulos.
\definecolor{cNucleo}{RGB}{36,122,82}
\definecolor{cImpl}{RGB}{186,86,30}
\definecolor{cEje}{RGB}{26,58,102}
\definecolor{cGris}{RGB}{120,116,110}

\tikzset{
  % los rótulos de la figura no se dividen con guion
  every node/.append style={execute at begin node={\hyphenpenalty=10000}},
  piloto/.style={circle, fill=#1, inner sep=0pt, minimum size=3.6mm},
  leyenda/.style={font=\small, anchor=west, text width=7.2cm, align=flush left},
  cifra/.style={font=\large\bfseries, text=#1, anchor=east},
}

\begin{document}

% --- Portada -----------------------------------------------------------------
\begin{titlepage}
\centering
\vspace*{3.2cm}
{\large\scshape Gestión del Desarrollo Tecnológico:\\
Estrategias y Análisis de Portfolio\par}
\vspace{0.5cm}
{\normalsize Módulo 6. Foro 6.1: Experiencia con proyectos de I+D\\
y nuevos retos\par}
\vspace{2.6cm}
{\color{cEje}\rule{0.55\textwidth}{0.8pt}\par}
\vspace{0.7cm}
{\LARGE\bfseries El bench como laboratorio\par}
\vspace{0.5cm}
{\large Cómo el tiempo entre contratos se convirtió\\
en innovación que sostiene el negocio\par}
\vspace{0.7cm}
{\color{cEje}\rule{0.55\textwidth}{0.8pt}\par}
\vfill
{\large Carlos Luis Rojas Aragonés\par}
\vspace{0.3cm}
{\normalsize Chief Technology Officer Program\par}
\vspace{0.3cm}
{\normalsize 10 de octubre de 2026\par}
\vspace*{1.5cm}
\end{titlepage}

% --- Cuerpo ------------------------------------------------------------------
\section*{1. El bench: holgura que se puede invertir}

En el desarrollo de software bajo el modelo de \emph{staff augmentation}
existe lo que informalmente llamamos el \enquote{bench}: ingenieros que
terminan un proyecto o contrato con un cliente y esperan el inicio del
siguiente. En lugar de tratarlos como costo ocioso, los usamos para formar un
departamento de innovación. \textcite{nohria1996} encontraron que cierta
holgura organizacional favorece la innovación, siempre que no sea excesiva;
el bench es exactamente esa holgura.

\section*{2. ¿En qué tipo de proyectos he participado?}

Sobre todo en pilotos de I+D aplicada. No inventamos tecnologías nuevas:
combinamos tecnologías ya estables y las pusimos a trabajar juntas para
resolver problemas reales.

\section*{3. ¿Cuál fue el resultado?}

De al menos 50 pilotos, 3 se implementaron y 2 forman hoy parte del núcleo de
la compañía, porque mantienen la operación funcionando de forma casi
automatizada (Figura~\ref{fig:pilotos}). Parece poco, pero
\textcite{stevens1997} estimaron que hacen falta unos 125 proyectos pequeños
para obtener un solo éxito comercial. El embudo es normal; lo importante es
que los ganadores paguen con creces los intentos que no prosperaron.

\begin{figure}[H]
\centering
\begin{tikzpicture}
  % Cuadrícula de 50 pilotos: 10 columnas por 5 filas
  \foreach \i in {0,...,49} {
    \pgfmathtruncatemacro{\col}{mod(\i,10)}
    \pgfmathtruncatemacro{\fila}{\i/10}
    \ifnum\i<2
      \def\tono{cNucleo}
    \else\ifnum\i<3
      \def\tono{cImpl}
    \else
      \def\tono{cGris!45}
    \fi\fi
    \node[piloto=\tono] at (0.55*\col, -0.55*\fila) {};
  }

  % Leyenda con las cifras
  \node[cifra=cNucleo] at (7.25,0)     {2};
  \node[leyenda] at (7.3,0)     {en el núcleo del negocio, operando hoy};
  \node[cifra=cImpl] at (7.25,-1.1)  {1};
  \node[leyenda] at (7.3,-1.1)  {implementado fuera del núcleo};
  \node[cifra=cGris] at (7.25,-2.2)  {47};
  \node[leyenda] at (7.3,-2.2)  {pilotos que no avanzaron, pero dejaron
    aprendizaje y práctica al equipo};

  \node[piloto=cNucleo] at (6.3,0)    {};
  \node[piloto=cImpl]   at (6.3,-1.1) {};
  \node[piloto=cGris!45] at (6.3,-2.2) {};
\end{tikzpicture}
\caption{Destino de los pilotos de I+D desarrollados con ingenieros del bench.
Cada círculo es un piloto: 3 de 50 llegaron a producción (6\,\%), una tasa
favorable frente a la referencia de un éxito por cada 125 proyectos pequeños
\parencite{stevens1997}.}
\label{fig:pilotos}
\end{figure}

\section*{4. El problema más importante: la financiación inestable}

De las opciones planteadas elijo la financiación inestable. El bench es
intermitente por naturaleza: cuando llega un contrato, el ingeniero vuelve a
facturar y el piloto se congela; además, cada hora de I+D compite con horas
facturables. Por eso cada proyecto debe nacer con una promesa atractiva y
alineada con la estrategia de la compañía, y pasar por puertas de decisión
explícitas, como en el modelo Stage-Gate de \textcite{cooper1990}, para
continuar, pausar o cerrar a tiempo.

\printbibliography[title={Referencias}]

\end{document}
